\section{Positive rows and distinct suffix words}\label{sec:positive}
Fix $r\ge1$. Start with the increasing seed $0,1,\ldots,r-1$. Append $r$ nonempty increasing blocks $B_1,\ldots,B_r$, of lengths $\ell_j\ge1$, making every suffix entry distinct across all blocks. Put $p_j=\ell_1+\cdots+\ell_j$ and $p_0=0$. For block $j$, select an $\ell_j$-element subset of
\begin{equation}\label{eq:lowerpool}
 \{0,1,\ldots,r+p_j-j-1\}\setminus
       \{\text{entries in earlier suffix blocks}\},
\end{equation}
and list it in increasing order. The ceiling increases by $\ell_j-1\ge0$ from the preceding block, so all earlier suffix entries belong to the current interval. Its candidate count is exactly
\begin{equation}\label{eq:rowchoices}
 r+\ell_j-j=\ell_j+(r-j+1)-1.
\end{equation}
Thus a row of width $k=r-j+1$ and total $\ell_j$ has precisely the right stars-and-bars count to parametrize the block.

\subsection{An explicit stars and bars rank bijection}
The correspondence is an actual reversible map, not merely equality of two counts. Let $y_1,\ldots,y_k\ge0$ sum to $\ell$, and set
\[
 b_i=y_1+\cdots+y_i+i\qquad(1\le i\le k-1).
\]
The $b_i$ are the increasing bar positions in $\{1,\ldots,\ell+k-1\}$. Take the complement of this bar set, an $\ell$-element set of star positions. Order the pool \eqref{eq:lowerpool} increasingly; its size is $\ell+k-1$ by \eqref{eq:rowchoices}. Select exactly the pool entries with those star ranks.

Conversely, a block and its decoded earlier suffix determine the pool and the ranks of the selected entries. The complementary ranks recover $b_1<\cdots<b_{k-1}$, and then
\[
 y_1=b_1-1,\quad y_i=b_i-b_{i-1}-1\ (2\le i<k),\quad
 y_k=\ell+k-1-b_{k-1}.
\]
For $k=1$, there are no bars and the unique row is $(\ell)$. This defines a bijection for every row and fixes the encoding unambiguously.

\begin{lemma}\label{lem:legalpositive}
Every word constructed above belongs to $\Av(000,100)$.
\end{lemma}
\begin{proof}
Before block $j$, the seed and internal block ascents alone guarantee at least
\[
 A_0=(r-1)+\sum_{h<j}(\ell_h-1)=r+p_{j-1}-j
\]
ascents. We deliberately ignore any ascents between blocks. If the selected entries of block $j$ are $x_1<\cdots<x_{\ell_j}$, their ceiling and their remaining successors imply
\[
 x_t\le r+p_j-j-1-(\ell_j-t)=A_0+t-1.
\]
For $t=1$, the bound is $x_1\le A_0$ and is legal. For $t\ge2$, the prefix before that entry has at least $A_0+t-2$ ascents, so \eqref{eq:ascent} holds. This remains true even if a block boundary is an ascent and two chosen blocks merge into one maximal increasing run.

Only seed labels can repeat, and each repeats at most once in the suffix. Thus $000$ is absent. For a seed label $v$, its first copy in the increasing seed precedes every larger label in the whole word. A larger entry cannot precede both copies of $v$, so $100$ is absent as well.
\end{proof}

\begin{lemma}\label{lem:blockfiber}
At a fixed suffix length $p\ge r$, the map from positive-row triangular arrays to the resulting words has fibers of size at most $\binom{p-1}{r-1}$.
\end{lemma}
\begin{proof}
Decorate a resulting word with its positive length composition $(\ell_1,\ldots,\ell_r)$. The seed length is fixed. Cutting its suffix at those lengths recovers all blocks, and the rank bijection recovers every row. Hence the decorated map is injective. There are exactly $\binom{p-1}{r-1}$ positive compositions of $p$ into $r$ parts.
\end{proof}
No claim is made that blocks can be recovered as maximal runs. For instance, an ascending boundary can disappear in the undecorated word. The explicit decoration and fiber bound account for this ambiguity.
